\chapter{Guion de animación}
\label{anx:guion_anim}
Este anexo presenta un guion técnico y su respectivo guion gráfico (\textit{storyboard}), 
los cuales sirven de base para la producción de una pieza audiovisual de carácter didáctico 
generada a partir de las herramientas de extracción y visualización tridimensional de datos intermedios de Tiny YOLOv2.
\medskip

El material está dirigido a estudiantes de programas académicos afines a las ciencias de la 
computación que cuenten con conocimientos fundamentales sobre aprendizaje automático y 
redes neuronales densas. Su objetivo principal consiste en ilustrar los conceptos nucleares 
de las redes neuronales convolucionales, tales como:
\begin{itemize}
    \item La operación formal de convolución bidimensional.
    \item La arquitectura jerárquica de la red (capas convolucionales y de reducción espacial o \textit{MaxPool2d}).
    \item Las divergencias estructurales respecto a redes densamente conectadas.
    \item El campo receptivo y la evolución representacional entre capas tempranas y profundas.
    \item Las diferencias funcionales entre tareas de clasificación y detección de objetos.
\end{itemize}

Dado su enfoque pedagógico, la propuesta se apoya en composiciones visuales limpias, transiciones fluidas y 
fondos de contraste oscuro, facilitando la identificación inmediata de los elementos volumétricos analizados. 
De forma complementaria, el componente visual se sincroniza con una locución en voz en \textit{off} para guiar 
la exposición, con una duración estimada de entre $1$ y $1.5$ minutos.

\section{Guión y Storyboard}
\centering
\small
\begin{xltabular}{\textwidth}{m{4.5cm} X X}
        % --- Encabezado en la primera página ---
    \caption{Guion técnico y correspondencia gráfica de la animación didáctica.} \label{tab:guion_tecnico_animacion} \\
    \toprule
    \textbf{Storyboard} & \textbf{Acción / Visual} & \textbf{Locución (Audio)} \\
    \midrule
    \endhead

    % --- Encabezado repetido en las páginas siguientes ---
    \toprule
    \textbf{Storyboard} & \textbf{Acción / Visual} & \textbf{Locución (Audio)} \\
    \midrule
    \endhead

    % --- Pie de tabla en páginas intermedias (indica que continúa) ---
    \midrule
    \multicolumn{3}{r}{\footnotesize\textit{Continúa en la siguiente página\dots}} \\
    \endfoot

    % --- Pie de tabla definitivo en la última página ---
    \bottomrule
    \endlastfoot

    \midrule
        \includegraphics[width=\linewidth, height=2.6cm, keepaspectratio]{storyboard/1.png} 
        & Vista frontal de la imagen de entrada. Se resalta el sujeto de la escena y se proyecta su caja 
        delimitadora de predicción. Transición hacia la vista tridimensional mientras se extruye la arquitectura de la red. 
        & Las redes neuronales convolucionales son arquitecturas especializadas, diseñadas para el análisis, 
        la clasificación y la detección de objetos en imágenes.\\
    \midrule
        \includegraphics[width=\linewidth, height=2.6cm, keepaspectratio]{storyboard/2.png}
        & Vista tridimensional en rotación continua del modelo Tiny YOLOv2.
        & Aunque comparten principios fundamentales con los perceptrones multicapa tradicionales, su principal 
        ventaja radica en la integración de capas especializadas capaces de procesar tensores volumétricos mediante 
        operaciones de convolución.\\
    \midrule
        \includegraphics[width=\linewidth, height=2.6cm, keepaspectratio]{storyboard/3.png}
        & Aparece una matriz de datos bidimensional sobre la cual se desplaza una ventana de cálculo. Se incorporan 
        dinámicamente elementos y filas para ilustrar el recorrido de la convolución discreta en dos dimensiones.
        & La convolución puede interpretarse como una ventana deslizante sobre un arreglo estructurado de datos, 
        como la cuadrícula de píxeles de una imagen digital.\\
    \midrule
        \includegraphics[width=\linewidth, height=2.6cm, keepaspectratio]{storyboard/4.png}
        & Alejamiento de cámara (\textit{zoom-out}): las celdas del arreglo reducen su escala relativa hasta 
        transformarse en los píxeles de una imagen, revelando finalmente la imagen completa.
        & Esta ventana, conocida como núcleo o \textit{kernel}, recorre la superficie realizando sumas ponderadas 
        con los valores del vecindario local, sintetizando un único valor representativo en cada posición espacial.\\
    \midrule
        \includegraphics[width=\linewidth, height=2.6cm, keepaspectratio]{storyboard/4.png}
        & El kernel transiciona a una suma ponderada con los valores en pantalla.
        & Una suma ponderada realiza una combinación lineal de los valores de una vecindad local, donde cada elemento aporta a la salida 
        proporcionalmente a un peso definido por el núcleo.\\
    \midrule
        \includegraphics[width=\linewidth, height=2.6cm, keepaspectratio]{storyboard/5.png}
        & Composición frontal con la imagen de entrada, el núcleo de convolución y el mapa de características resultante. 
        El núcleo y la salida rotan secuencialmente a modo de carrusel.
        & En el procesamiento digital clásico, estas transformaciones se emplean para suavizar señales, atenuar ruido o 
        resaltar características morfológicas clave, como los bordes de una figura. A estas transformaciones se les conoce como filtros.\\
    \midrule
        \includegraphics[width=\linewidth, height=2.6cm, keepaspectratio]{storyboard/6.png}
        & La imagen de entrada y el mapa de salida se segmentan en regiones independientes, cada una evaluada por su 
        respectivo núcleo de convolución en paralelo.
        & Además, al constituir evaluaciones locales e independientes por región, la convolución es una operación 
        intrínsecamente paralelizable y de alta eficiencia computacional.\\
    \midrule
        \includegraphics[width=\linewidth, height=2.6cm, keepaspectratio]{storyboard/6.png}
        & Difuminado a modelo tridimensional de un GPU.
        & Esto permite que las convoluciones se beneficien de la aceleración por hardware especializado, como las unidades de procesamiento gráfico (\textit{GPU}), 
        que están diseñadas para ejecutar operaciones gráficas de forma eficiente.\\
    \midrule
        \includegraphics[width=\linewidth, height=2.6cm, keepaspectratio]{storyboard/7.png}
        & La imagen de salida transiciona al espacio tridimensional y se descompone en dos retículas de vóxeles que 
        representan los elementos de entrada y salida. Se trazan aristas de interconexión densa entre cada par de elementos.
        & Sin embargo, intentar procesar imágenes de alta resolución mediante capas totalmente conectadas presenta un obstáculo 
        crítico: vincular cada par de neuronas entre capas sucesivas incrementa drásticamente el número de parámetros y la complejidad computacional.\\
    \midrule
        \includegraphics[width=\linewidth, height=2.6cm, keepaspectratio]{storyboard/8.png}
        & Acercamiento de cámara (\textit{zoom-in}) hacia el volumen de la capa convolucional, enfocando el 
        prisma alámbrico del núcleo.
        & Las redes convolucionales resuelven este problema mediante las capas convolucionales. Los pesos son compartidos mediante los núcleos convolucionales: un
        mismo conjunto de parámetros se reutiliza a lo largo de toda la superficie de la capa, aprendiendo a 
        identificar patrones espaciales específicos con independencia de su posición en el encuadre.\\
    \midrule
        \includegraphics[width=\linewidth, height=2.6cm, keepaspectratio]{storyboard/8.png}
        & Se muestra la operación de convolución entre dos volumenes, mostrando como se apilan los mapas de activación de salida para formar el volumen de activación.
        & Así transforman tensores volumétricos mediante la aplicación de un conjunto de filtros tridimensionales.
        Cada filtro realiza un barrido sobre el tensor de entrada, generando un mapa de características o canal por cada uno, que representa la presencia del patrón.
        Y los canales se apilan para formar un nuevo tensor volumétrico conteniendo las activaciones de la capa.\\
    \midrule
        \includegraphics[width=\linewidth, height=2.6cm, keepaspectratio]{storyboard/9.png}
        & Se muestran varios núcleos con valores de pesos que cambian.
        & Los pesos de los núcleos utilizados por estos filtros se ajustan durante el entrenamiento para optimizar la detección de patrones específicos.\\
    \midrule
        \includegraphics[width=\linewidth, height=2.6cm, keepaspectratio]{storyboard/9.png}
        & Visualización de la operación de reducción (\textit{pooling}). Se delimita una vecindad de cuatro píxeles 
        adyacentes y se descartan tres; el proceso se repite sobre toda la cuadrícula y las activaciones retenidas se compactan en una nueva retícula.
        & Otra operación fundamental en estas arquitecturas es el agrupamiento o \textit{pooling}, cuyo propósito es 
        reducir la dimensionalidad espacial de las representaciones intermedias. Su variante más común, el 
        \textit{max-pooling}, reduce las dimensiones transversales a la mitad reteniendo únicamente la activación máxima de cada vecindad.\\
    \midrule
        \includegraphics[width=\linewidth, height=2.6cm, keepaspectratio]{storyboard/10.png}
        & Vista frontal del grafo. Las capas intermedias se ocultan progresivamente para evidenciar la reducción de 
        resolución espacial. Se proyecta un indicador cuadrangular sobre un vóxel profundo; luego, las capas 
        reaparecen en orden inverso mientras el recuadro se expande progresivamente hacia la entrada para ilustrar la cobertura del campo receptivo.
        & Este mecanismo de compresión alivia la carga de cómputo en las etapas subsiguientes y expande 
        progresivamente el campo receptivo de la red.\\
    \midrule
        \includegraphics[width=\linewidth, height=2.6cm, keepaspectratio]{storyboard/10.png}
        & Vista lateral que muestra el campo receptivo de varias capas, demostrando cómo se amplía a medida que se profundiza en la red.
        & El campo receptivo es el área de la imagen de entrada que influye en la activación de una neurona específica. Capas más profundas tienen campos receptivos más grandes,
        permitiendo abarcar regiones más amplias de la imagen.\\
    \midrule
        \includegraphics[width=\linewidth, height=2.6cm, keepaspectratio]{storyboard/11.png}
        & Despliegue comparativo en paralelo de los mapas de activación correspondientes a distintas profundidades de la red.
        & De este modo, las capas iniciales responden a rasgos locales primitivos, como bordes, transiciones de 
        contraste y texturas elementales, mientras que las capas más profundas sintetizan estas respuestas para 
        reconocer estructuras abstractas y de mayor escala.\\
    \midrule
        \includegraphics[width=\linewidth, height=2.6cm, keepaspectratio]{storyboard/12.png}
        & Retorno a la perspectiva global de la arquitectura tridimensional; se iluminan sucesivamente 
        los bloques de capas a medida que se enuncian en la locución.
        & En términos generales, una red convolucional intercala secuencias de convolución, normalización 
        por lotes, activaciones no lineales y agrupamiento espacial para estructurar una jerarquía progresiva de representaciones visuales.\\
    \midrule
        \includegraphics[width=\linewidth, height=2.6cm, keepaspectratio]{storyboard/13.png}
        & Visualización del plano de salida junto a un listado de clases ordenadas descendentemente según 
        su probabilidad predicha.
        & La topología final del modelo se adapta a la tarea asignada: en problemas de clasificación pura, 
        la red genera una distribución categórica de probabilidades que asocia la totalidad de la imagen a una única etiqueta representativa.\\
    \midrule
        \includegraphics[width=\linewidth, height=2.6cm, keepaspectratio]{storyboard/15.png}
        & Presentación de la matriz estructurada de predicciones. Los vectores se trasladan hacia el plano 
        de salida y se resaltan individualmente, proyectando sus respectivas cajas y etiquetas tridimensionales en la escena.
        & En detección de objetos, arquitecturas más avanzadas identifican múltiples instancias concurrentes, generando 
        como salida un conjunto de vectores donde cada registro codifica: las coordenadas espaciales y dimensiones 
        de la caja delimitadora, la puntuación de objetualidad, la distribución de 
        probabilidad condicional, y el identificador de clase.\\
\end{xltabular}
